\documentclass{article}
\usepackage{hyperref}
\usepackage{Style}

\nocite{*} % Comentar si quiero citar
%\addbibresource{bibliografia.bib} % Quitar el comentado si quiero usar bibliografia

\begin{document}

\begin{minipage}{2.5cm}
    \includegraphics[width=2cm]{imagen_puc.jpg}
\end{minipage}
\begin{minipage}{12cm}
    {\sc Pontificia Universidad Católica de Chile\\
    Facultad de Matemáticas\\
    Departamento de Matemática\\
    Profesor: Manuel Cabezas -- Ayudante: Benjamín Mateluna}
\end{minipage}
\vspace{2ex}

{\centerline{\bf Introducción a la Geometría - MAT1304}
\centerline{\bf Ayudantía 26 - Repaso Geometría Euclideana}}
\centerline{\bf 23 de junio de 2026}

\vspace{1cm}
\noindent\textbf{Problema 1.} Considere el segmento $\overline{AB}$. Sean $C$ y $D$ puntos tales
que $\overline{AC}$ y $\overline{BD}$ son perpendiculares a $\overline{AB}$ y además 
$\overline{AC}=\overline{BD}$. Pruebe que $\overline{CD}$ es paralelo a $\overline{AB}$.

\vspace{5mm}
\noindent\textbf{Problema 2.} Dada una recta $L$ considere los puntos $M$ y $N$ al mismo lado de 
$L$. Construya un punto $C$ en $L$ tal que el ángulo que forma $L$ con $\overline{MC}$ es 
congruente con el ángulo determinado por $L$ con $\overline{NC}$.

\vspace{5mm}
\noindent\textbf{Problema 3.} Sean $\mathcal{C}_{1}$ y $\mathcal{C}_{2}$ dos circunferencias con
centros $O_{1}$ y $O_{2}$ y radios $r_{1}>r_{2}>0$, cada una externa a la otra. Construir con 
regla y compás una recta que sea tangente a ambas circunferencias.

\vspace{5mm}
\noindent\textbf{Problema 4.} Dada una circunferencia de centro $O$, construir un triángulo 
isósceles que lo tenga como íncirculo, es decir, cuyos lados sean tangentes a la circunferencia.

\vspace{5mm}
\noindent\textbf{Problema 5.} Demuestre la identidad
\begin{equation*}
    cot\left(\frac{x}{2}\right)=\frac{1+cos(x)}{sen(x)}
\end{equation*}

\vspace{1cm}
\noindent\textbf{\underline{Ejercicios Propuestos:}}

\vspace{5mm}
\noindent\textbf{Extra 1.} Considerar un triángulo $\triangle ABC$ cualquiera y $\overline{BD}$ la 
bisectriz desde $B$. Sea $E\in\overline{BC}$ tal que $\overline{DE}$ es paralelo a 
$\overline{AB}$. Asumir que $\overline{DE}=3$ y $\overline{BC}=3\overline{AB}$. Encontrar 
$\overline{BC}$.

\vspace{5mm}
\noindent\textbf{Extra 2.} Mostrar que $cos(\frac{\pi}{9})cos(\frac{2\pi}{9})cos(\frac{4\pi}{9})
=\frac{1}{8}$.

%\printbibliography % Quitar el comentado si quiero usar bibliografia

\end{document}